\chapter{Asthma in Children (6 to 11 Years)}
\index{Asthma in Children (6 to 11 Years)}

\section*{Definition and Overview}
Asthma is a chronic inflammatory condition of the airways that makes them swollen, narrowed, and overly sensitive. In a child aged 6 to 11 years, the muscles around the breathing tubes can tighten, the airway lining can swell, and excess mucus can be produced, all of which narrow the airways and make breathing difficult. Between episodes many children feel completely well, and the condition is characterised by recurring, variable symptoms. Asthma at this age is usually diagnosed on the basis of repeated symptoms of wheeze, cough, chest tightness, or breathlessness, together with a documented response to treatment. With the right medicines, correct inhaler technique, and an individualised asthma action plan, most school-age children can exercise, play sport, and attend school without limitation.

\section*{Epidemiology}
Asthma is one of the most common chronic conditions of childhood worldwide, affecting about one in ten children in many populations. It is slightly more common in boys than in girls before puberty, although this difference tends to even out in adolescence. Asthma frequently runs in families and is more common in children with eczema, hay fever, or food allergies. Symptoms are often worse during respiratory infections, pollen seasons, cold weather, and exposure to tobacco smoke. Many children with asthma at this age improve as they grow older, although a substantial proportion continue to have symptoms into adult life.

\section*{Aetiology and Pathophysiology}
Asthma results from a combination of genetic susceptibility and environmental triggers. The exact cause is not fully understood, but a family history of allergy or asthma is a strong risk factor.

\begin{itemize}
    \item \textbf{Inflammation:} the airway lining is chronically inflamed with white blood cells, such as eosinophils, even when the child feels well
    \item \textbf{Bronchoconstriction:} the muscle layer around the airways contracts in response to triggers, narrowing the tubes
    \item \textbf{Mucus hypersecretion:} glands in the airway wall produce increased sticky mucus, which can plug the small airways
    \item \textbf{Airway hyper-responsiveness:} the airways overreact to stimuli that would not trouble a child without asthma
    \item \textbf{Common triggers:} viral colds, exercise, cold air, house dust mites, pollens, pet dander, mould, tobacco smoke, and strong emotions
\end{itemize}

\section*{Clinical Features}
Symptoms tend to come and go, and are often worse at night, in the early morning, or after exercise. Each child has their own pattern of symptoms and triggers.

\begin{itemize}
    \item A whistling sound on breathing out (wheeze)
    \item A dry cough, often worse at night, with exercise, or when laughing
    \item Shortness of breath, or a feeling of being chesty
    \item Tightness or discomfort in the chest
    \item Tiredness, poor sleep, and irritability during flare-ups
    \item Symptoms that come on with a known trigger, such as a cold or exercise
\end{itemize}

\section*{Differential Diagnosis}
Not every child who wheezes has asthma, and the diagnosis can be uncertain in this age group.

\begin{itemize}
    \item \textbf{Viral wheeze:} wheeze occurring only with viral colds, which is common in younger children
    \item \textbf{Bronchiolitis:} a viral lower respiratory infection, usually in children under two years
    \item \textbf{Foreign body aspiration:} sudden wheeze or cough after choking on a small object or food
    \item \textbf{Gastro-oesophageal reflux:} acid washing back into the airways can cause cough and wheeze
    \item \textbf{Allergic rhinitis and sinusitis:} post-nasal drip causing a chronic cough
    \item \textbf{Vocal cord dysfunction:} noisy breathing that mimics asthma but is not due to bronchoconstriction
    \item \textbf{Cystic fibrosis, congenital heart disease, and other lung problems:} rare causes of persistent symptoms
\end{itemize}

\section*{When to Seek Medical Care}
See a general practitioner if the child is wheezing, coughing, or short of breath for the first time, needs the reliever more often than every four hours, wakes at night with symptoms, or if symptoms are interfering with school or sport. A review is also needed after any flare-up or emergency attendance.

\textbf{Contact your local emergency services for:}
\begin{itemize}
    \item Difficulty breathing or severe shortness of breath
    \item Chest in-drawing, where the skin is sucked in between the ribs, or grunting with breathing
    \item Struggling to speak in full sentences
    \item Lips or face turning blue
    \item The reliever not helping after several puffs
\end{itemize}

\section*{Diagnosis and Investigations}
In this age group the diagnosis is usually made clinically and confirmed over time by the pattern of symptoms and the response to treatment.

\begin{itemize}
    \item \textbf{Detailed history:} symptoms, triggers, night waking, family history of allergy, and response to relievers
    \item \textbf{Physical examination:} listening to the chest, and checking for eczema and signs of allergic rhinitis
    \item \textbf{Spirometry:} measurement of lung function by blowing into a machine, achievable in most children from about age 6
    \item \textbf{Peak flow monitoring:} a handheld meter that measures how fast air can be blown out
    \item \textbf{Reversibility testing:} repeating lung function after a reliever to see whether the airways open up
    \item \textbf{Allergy testing:} skin-prick or blood tests for allergens if triggers are unclear
    \item \textbf{Trial of treatment:} a trial of a preventer inhaler with review of symptoms over several weeks
\end{itemize}

\section*{Management}
Asthma at this age is managed with a written asthma action plan agreed with the doctor, regular review, and medicines taken with correct technique.

\begin{itemize}
    \item \textbf{Reliever (such as salbutamol):} used when needed to relieve symptoms quickly; frequent use signals poor control and the need for review
    \item \textbf{Preventer inhaler:} a daily inhaled corticosteroid to reduce airway inflammation and prevent flare-ups
    \item \textbf{Inhaler technique:} spacer devices help children take the medicine into the lungs correctly
    \item \textbf{Combination or add-on medicines:} considered if symptoms are not controlled with a preventer alone
    \item \textbf{Flare-up management:} follow the action plan, increase reliever use, and see a doctor promptly
    \item \textbf{Exercise:} with good control, children should be able to join in all activities, and sport should remain part of normal life
    \item \textbf{School support:} give the school a copy of the action plan and ensure reliever and spacer are accessible
    \item \textbf{Regular review:} at least every 6 to 12 months, and more often if control is poor
\end{itemize}

\section*{Complications}
Poorly controlled asthma can affect every part of a child's life and, rarely, can be dangerous.

\begin{itemize}
    \item Severe asthma flare-ups requiring hospital admission or intensive care
    \item Missed school days, and interrupted sleep and play
    \item Difficulty keeping up with sport and physical activity
    \item Slower growth velocity with long-term severe symptoms or frequent oral corticosteroids
    \item Anxiety and low mood, partly related to the unpredictability of attacks
    \item Side effects of medicines, mainly with high-dose inhaled steroids or repeated courses of oral steroids
    \item In rare cases, a severe attack can be life-threatening
\end{itemize}

\section*{Prognosis and Outcomes}
Asthma is a long-term condition, but the outlook for school-age children is generally good. With effective treatment, most children have no limitation of normal activities, and many appear to grow out of troublesome symptoms during adolescence. However, symptoms may return in adult life, and the underlying airway hyper-responsiveness often persists. Children whose symptoms began in the preschool years, who have strong allergies, or who are exposed to tobacco smoke are more likely to have persistent asthma. Regular review and adherence to the action plan are the strongest predictors of good control.

\section*{Prevention and Self-Care}
Asthma cannot be cured or reliably prevented, but flare-ups and their impact can be greatly reduced by good daily care.

\begin{itemize}
    \item Take preventer medicines exactly as prescribed, even when the child feels well
    \item Keep the written asthma action plan updated and follow it during flare-ups
    \item Reduce exposure to tobacco smoke, and never smoke around children
    \item Keep allergens such as house dust mites, mould, and animal dander under control where they are triggers
    \item Encourage normal physical activity, using a reliever before exercise if advised
    \item Treat hay fever, because it often makes asthma worse
    \item Have the influenza vaccination each year if recommended
    \item Bring the reliever and spacer to school, sport, and holidays
    \item Attend regular asthma reviews, even when the child is well
\end{itemize}

\section*{Sources and Further Reading}
The following organisations provide reliable, evidence-based information on asthma in children for parents and healthcare students.

\begin{itemize}
    \item Global Initiative for Asthma (GINA). \textit{Pocket guide for asthma management and prevention in children.} https://ginasthma.org/
    \item NHS. \textit{Asthma in children: symptoms, treatment, and support.} https://www.nhs.uk/conditions/asthma/
    \item Asthma + Lung UK. \textit{Childhood asthma information for parents.} https://www.asthmaandlung.org.uk/
    \item MedlinePlus. \textit{Asthma in children.} https://medlineplus.gov/asthmainchildren.html
    \item Centers for Disease Control and Prevention (CDC). \textit{Asthma in children.} https://www.cdc.gov/asthma/
\end{itemize}
